\section{Introduction}

\label{2305:sec:intro}
In this manuscript we are interested in the supervised task of learning the following target function:
\begin{align}
\label{2305:eq:def:teacher}
    y = \sigma_{\star}\left(w_{\star}^{\top}x\right) + \sqrt{\Delta}z, && x\sim\mathcal{N}(0,\sfrac{1}{d} I_{d}), \qquad z\sim\mathcal{N}(0,1).
\end{align}
This target belongs to a general class of models known as \emph{single-index} or \emph{generalized linear} models, where the labels depend on the covariates $x\in\mathbb{R}^{d}$ only through its projection on a fixed direction $w_{\star}\in\mathbb{R}^{d}$, followed by a non-linear real-valued function $\sigma_{\star}:\mathbb{R}\to\mathbb{R}$. Despite its apparent simplicity, the study of single-index models with Gaussian i.i.d. covariates has been at the core of the recent theoretical progress in machine learning \citep{tan2019online,arous2021online,bietti2022learning,damian2022neural,berthier2023learning}. Note that if $\sigma_{\star}(x)=x$, learning could be achieved efficiently by linear methods such as ridge regression or PCA. The presence of a non-linearity $\sigma_{\star}$ makes it perhaps the simplest target for which a non-linear hypothesis class is required. Moreover, the choice of isotropic covariates $x\sim\mathcal{N}(0,\sfrac{1}{d}I_{d})$ means that all structure in the target is in learning a representation of $\sigma$ and a particular direction $w_{\star}$. The popularity of this model is that different separation results can be shown in the high-dimensional limit where $d\gg 1$:
\begin{itemize}
    \item In the well-specified setting where we fit an isotropic single-index model with the same hypothesis class (i.e. $f_{\theta}(x) = \sigma_{\star}(w^{\top}x)$), it has been shown that the sample complexity of one-pass SGD\footnote{Which we recall the reader is equivalent to the convergence rate.}  is determined by the first non-zero Hermite coefficient of the target $\sigma_{\star}$, also known as the \emph{information exponent} \citep{arous2021online}. Problems with non-zero first Hermite have information exponent $\ell=1$, and $w_{\star}$ can be learned at linear sample complexity $n=O(d)$. Instead, problems with zero first and non-zero second Hermite coefficient have information exponent $\ell=2$, requiring instead $n = O(d\log{d})$ samples \citep{tan2019online,arous2021online}. Note that this is a factor $\log{d}$ higher than the Bayes-optimal sample complexity of $n=O(d)$ \citep{barbier2019optimal,maillard2020phase}, hence defining a class of \emph{hard} problems for SGD. 
    \item For fully-connected two-layer neural networks $f_{\theta}(x) = a^{\top}\sigma(Wx)$, several results are known under different assumptions. For fixed first layer weights $W\in\mathbb{R}^{p\times d}$ (a.k.a. random features model) and large enough width $p$, learning the $k$-th order Hermite coefficient $\sigma_{\star}$ requires $n=O(d^{k})$ samples \citep{mei2022generalizationerror}, implying a sample complexity of $n=O(d^2)$ for a quadratic problem, e.g. $\sigma_{\star}(x)=x^2$. Recently, it was shown that wide networks ($p\to\infty$) can achieve the well-specified sample complexity of $n=O(d)$ under one-pass SGD, \emph{provided that all Hermite coefficients of both $\sigma_{\star}, \sigma$ are non-zero} \citep{berthier2023learning}. In particular, this assumption cover only problems with information exponent $\ell=1$, excluding hard cases such as quadratic problems. Finally, for $\sigma(x) = \sigma_{\star}(x) = x^2$, \cite{saraomannelli2020optimization} has shown that for $p$  large enough, full-batch gradient flow achieves sample complexity $n=2d$, although at a running time of $t=O(\log{d})$. 
\end{itemize}
With the exception of \cite{berthier2023learning}, the works mentioned above cover the scaling of the sample complexity in the high-dimensional limit. Our goal is, instead, to derive sharp results for the sample complexity of learning \eqref{2305:eq:def:teacher} with a fully-connected two-layer neural network in the challenging case where $\sigma_{\star}$ has a vanishing first Hermite coefficient.
As discussed above, this case violates the ``standard learning scenario'' of \cite{berthier2023learning}, and can be seen as a proxy for hard learning problems for descent-based algorithms. For concreteness, in the following we focus on the purely quadratic case:
\begin{align}
\label{2305:eq:def:teacher:quadratic}
 y = \left(w_{\star}^{\top}x\right)^2  + \sqrt{\Delta}z, && w_{\star}\in\mathbb{S}^{d-1}(\sqrt{d})
\end{align}
%The key challenge is that the quadratic target \eqref{2305:eq:def:teacher} has zero leading Hermite coefficient, violating the "standard learning scenario" of \cite{berthier2023learning}. It therefore can be seen as a proxy model for a hard learning task for descent-based algorithms. 
Learning the target \eqref{2305:eq:def:teacher:quadratic} consists of learning the non-linearity $\sigma_{\star}(x)=x^2$ and the direction $w_{\star}$. In this work, we focus our attention in the second part, considering the following architecture with squared-activation: 
\begin{align}
\label{2305:eq:student}
    f_{\Theta}(x) = \frac{1}{p}\sum\limits_{i=1}^{p}a_{i}(w_{i}^{\top}x)^2.
\end{align}
where $\Theta = (a,W)$ is the set of trainable weights, which are trained with one-pass stochastic gradient descent (SGD):
%Since our focus is on learning the direction $w_{\star}\in\mathbb{S}^{d-1}(\sqrt{d})$, we normalize the weights $w_{i}\in\mathbb{S}^{d-1}(\sqrt{d})$, 
\begin{align}
\label{2305:eq:sgd}
    % a_{i}^{\nu+1} &= a^{\nu}_{i} - \gamma_{a} \nabla_{a_{i}}\ell(y^{\nu}, f_{\Theta}(x^{\nu}))\notag\\
    % w_{i}^{\nu+1} &= {\rm Proj}_{\mathbb{S}^{d-1}(\sqrt{d})}\left(w_{i}^{\nu} - \gamma \nabla_{w_{i}}\ell(y^{\nu}, f_{\Theta}(x^{\nu}))\right),
    \Theta^{\nu+1} = \Theta^{\nu} -\gamma \nabla_{\Theta}\ell(y^{\nu},f_{\Theta^{\nu}}(x^{\nu}))
\end{align}
with square loss $\ell(y,x)=\sfrac{1}{2}(y-x)^2$ and initial condition $\Theta^{0} = (a^{0}, W^{0})$. Note that at each step $\nu$, the gradient is evaluated at a fresh pair of data $(x^{\nu}, y^{\nu})\in\mathbb{R}^{d+1}$ drawn from the model \eqref{2305:eq:def:teacher:quadratic}.
%and the projection in the sphere ensures the weights $w_{i}^{\nu}\in\mathbb{S}^{d-1}$ remain normalized throughout training. 
In particular, this implies that after $\nu\in[n]$ steps, the algorithm has seen $n$ data points. 

Learning in this problem is hard, and can be compared to finding a needle in a haystack. Indeed, with the exception of one direction that points towards $\pm \w_{\star}$, the population risk at (random) initialization is mostly flat. This slows down the dynamics, which takes a long time to establish a significant correlation with the signal - a scenario we refer to as \emph{escaping mediocrity}. 

At first, the particular case of purely quadratic activation might appear too specific. Indeed, as we will see later the population risk for this task has a global maximum at initialization and a degenerated set of global minima. The choice of more general $\sigma_{\star}$ and $\sigma$ with zero first Hermite but not necessarily zero higher-order coefficients might give rise to other critical points such as saddle-points, giving rise to a more complex SGD dynamics. However, since the focus of this work is on escaping mediocrity, our conclusions will hold, up to constants, to more general activations with information exponent \(\ell=2\).

\paragraph{Summary of results ---} Our main contributions in this manuscript are: 
\begin{itemize}
    \item We derive a deterministic set of ODEs providing an exact and analytically tractable description of the one-pass SGD dynamics in the high-dimensional limit $d\to\infty$, and characterize the leading order stochastic corrections to this limit. 
    \item We provide an analytical formula for the number of samples required for one-pass SGD to learn the phase retrieval target in high-dimensions at arbitrary network width. We show that overparameterization can only improve convergence by a constant factor for phase retrieval.
    \item Finally, we compute the leading order stochastic corrections to the exit time, and show that stochasticity does not help to escape the flat directions at initialization. This suggests that the deterministic descriptions is enough to fully capture the phenomenology of the dynamics in this problem.
\end{itemize}

All the codes used for numerical experiments are provided in this \href{https://github.com/IdePHICS/EscapingMediocrity}{GitHub repository}.

\paragraph{Further related work ---}
The investigation of a deterministic high-dimensional limit of one-pass SGD for two-layer neural networks draws back to the seminal works of \cite{saad1995exact,saad1995line,saad1996dynamics}, and was followed by a stream of works that span decades of research \citep{riegler1995line,copelli1995line,reents1998self,goldt2019dynamics,goldt2020modeling,goldt2022gaussian,veiga2022phase,arnaboldi2023high}. More recently, the stochastic corrections around fixed points of the dynamics have been investigated by \cite{arous2022high}. In a complementary research line, \cite{mei2018mean,chizat2018global,rotskoff2022trainability,ghorbani2019limitations,sirignano2020mean} have shown that an alternative deterministic description of SGD can be obtained in the infinite width-limit, a.k.a. mean-field regime. High-dimensional reductions of the mean-field equations have been studied by \cite{abbe2022merged,abbe2023sgd,hajjar2023symmetries,berthier2023learning}. Recently, \cite{veiga2022phase,arnaboldi2023high} has shown that these apparently different limits of one-pass SGD can be unified in a single description. 

There has been a recent surge of interest in studying how increasing degrees of complexity in the target function are incrementally learned by SGD \citep{abbe2022merged,abbe2023sgd,jacot2022saddle,boursier2022gradient,berthier2023learning}, with an emerging staircase picture where complexity is sequentially learned in different scenarios. This picture, however, is bound to classes of targets where SGD develops strong correlations with the target directions at initialization, a notion which was mathematically formalized by the so-called information exponent \(\ell\) by \cite{arous2021online}. Instead, targets for which the landscape at initialization is mostly flat (\(\ell\geq 2\)) are hard for SGD at high-dimensions, translating to very slow dynamics. This is precisely the case for the phase retrieval problem (\(\ell = 2\)), a classic inverse problem arising in many scientific areas, from X-ray crystallography to astronomical imaging \citep{jaganathan2016phase,dong2023phase}. Phase retrieval has been widely studied in the literature as a prototypical example of a hard inverse problem \citep{candes2013phase,candes2013phaselift,candes2015phase,mondelli2018fundamental,barbier2019optimal,aubin2020exact,maillard2020phase}, providing a simple yet challenging example of a non-convex optimization problem which is hard for descent-based algorithms \citep{sun2016geometric,tan2019online,tan2018phase,chen2019gradient,davis2020nonsmooth,saraomannelli2020optimization,saraomannelli2020complex,mignacco2021stochasticity}. 
%%%%%%%%%%%%%%%%%%%%%%%%%%%%%%%%%%%%%%%%%%%%%%%%%%%%%%%%%%%%%%%%%%%%%%%%%%%%%%%
